\section{Two-level relaxation by phonons and electrons}\label{app:D.42}

For the coupling $\tfrac12\sigma_z X$ of a two-level system to a bath coordinate $X$, \cref{eq:16.8-golden} needs only the spectral density $J_{\mathrm{s}}(\omega)$ of \cref{eq:C.20-J}. The notation is that of \cref{app:C.20,sec:16.8}: asymmetry $\Delta_{\mathrm{as}}$, tunneling amplitude $\Delta_0$, splitting $\Delta_{\mathrm{tl}}$, and $\sigma_z = |\mathrm{L}\rangle\langle\mathrm{L}| - |\mathrm{R}\rangle\langle\mathrm{R}|$.

\subsection*{Acoustic phonons}

Let a strain $u_{\alpha\beta}$ shift the asymmetry by $\delta\Delta_{\mathrm{as}} = 2\sum_{\alpha\beta}\gamma_{\alpha\beta}u_{\alpha\beta}$, so that the coupling is $\sigma_z\sum_{\alpha\beta}\gamma_{\alpha\beta}u_{\alpha\beta}$, with the deformation-potential tensor $\gamma_{\alpha\beta}$ of the defect. This is the convention of \cite[Eq.~(16)]{muller2019towards} and of the standard tunneling model, in which the relaxation rate takes the form \cref{eq:16.7-tau1}. The text of Meißner et al.\ states $\delta\varepsilon = \gamma u_{\mathrm{d}}$ for the energy shift in a strain component $u_{\mathrm{d}}$, but their one-phonon rate~\cite[Eq.~(7)]{meissner2021acoustic} holds for $\delta\varepsilon = 2\gamma u_{\mathrm{d}}$; with $\delta\Delta_{\mathrm{as}} = \gamma S$ it would be four times smaller. In an isotropic solid of density $\rho_T$ and volume $\Vol$ the displacement field at long wavelength is the continuum limit of \cref{sec:5.3},
\begin{equation*}
  \vect{u}(\vect{r}) = \sum_{\vect{Q}\nu}\frac{1}{(2\rho_T\Vol\omega_{\vect{Q}\nu})^{1/2}}\,\hat{\vect{e}}_{\vect{Q}\nu}\bigl(d_{\vect{Q}\nu}\ee^{\ii\vect{Q}\cdot\vect{r}} + d_{\vect{Q}\nu}^\dagger\ee^{-\ii\vect{Q}\cdot\vect{r}}\bigr) ,
  \qquad \omega_{\vect{Q}\nu} = c_\nu Q ,
\end{equation*}
with one longitudinal branch, $\hat{\vect{e}} = \hat{\vect{Q}}$, and two transverse ones. At the defect, placed at the origin, the strain $u_{\alpha\beta} = \tfrac12(\partial_\alpha u_\beta + \partial_\beta u_\alpha)$ of one mode is $\ii Q\,\tfrac12(\hat Q_\alpha\hat e_\beta + \hat Q_\beta\hat e_\alpha)(2\rho_T\Vol\omega_{\vect{Q}\nu})^{-1/2}(d_{\vect{Q}\nu} - d_{\vect{Q}\nu}^\dagger)$. The phase $\ii$ is removed by redefining $d_{\vect{Q}\nu} \to -\ii d_{\vect{Q}\nu}$, and the coupling takes the form $\tfrac12\sigma_z\lambda_{\vect{Q}\nu}(d_{\vect{Q}\nu} + d_{\vect{Q}\nu}^\dagger)$ with
\begin{equation*}
  \lambda_{\vect{Q}\nu} = \frac{2Q\,\gamma_{\vect{Q}\nu}}{(2\rho_T\Vol\omega_{\vect{Q}\nu})^{1/2}} , \qquad
  \gamma_{\vect{Q}\nu} = \sum_{\alpha\beta}\gamma_{\alpha\beta}\hat Q_\alpha\hat e_\beta ,
\end{equation*}
where the symmetry of $\gamma_{\alpha\beta}$ has been used. The squares $\gamma_{\vect{Q}\nu}^2$ depend on the direction of $\vect{Q}$ relative to the defect. Averaged over the orientations of the defects of a glass, or equivalently over $\hat{\vect{Q}}$ for one defect, they define $\gamma_{\mathrm{l}}^2 = \langle(\hat{\vect{Q}}\cdot\gamma\cdot\hat{\vect{Q}})^2\rangle$ and, for each transverse polarization, $\gamma_{\mathrm{t}}^2 = \langle(\hat{\vect{Q}}\cdot\gamma\cdot\hat{\vect{e}}_{\mathrm{t}})^2\rangle$. Then $\lambda_{\vect{Q}\nu}^2 = 2\gamma_\nu^2\omega/(\rho_T\Vol c_\nu^2)$ at $\omega = c_\nu Q$, and with $\sum_{\vect{Q}} \to \Vol\int\dd^3Q/(2\pi)^3$,
\begin{equation*}
  J_{\mathrm{s}}^{\mathrm{ph}}(\omega) = \pi\sum_\nu\int\frac{\Vol\,4\pi Q^2\dd Q}{(2\pi)^3}\,\frac{2\gamma_\nu^2\,c_\nu Q}{\rho_T\Vol c_\nu^2}\,\delta(\omega - c_\nu Q)
  = \sum_\nu\frac{\gamma_\nu^2\,\omega^3}{\pi\rho_Tc_\nu^5} ,
\end{equation*}
which with one longitudinal and two transverse branches is \cref{eq:16.8-Jph}. In \cref{eq:16.8-golden} it gives \cref{eq:16.7-tau1}. The form of Sahling et al., $[A\Delta_{\mathrm{tl}}\Delta_0^2\coth(\Delta_{\mathrm{tl}}/2T)]$ with $A = (8\pi^3/\rho h^4)(\gamma_{\mathrm{l}}^2/v_{\mathrm{l}}^5 + 2\gamma_{\mathrm{t}}^2/v_{\mathrm{t}}^5)$~\cite[Eqs.~(4), (5)]{sahling2010anomalous}, is the same with $h = 2\pi$ and $k_{\mathrm{B}} = 1$. The long-wavelength form needs $\Delta_{\mathrm{tl}}/c_\nu$ small compared with the inverse size of the defect, which at \SI{0.1}{meV} and \SI{4}{km.s^{-1}} is a wavelength of \SI{165}{nm} against a defect of a few \r{a}ngstr\"oms.

\subsection*{Conduction electrons in the Born limit}

Let the defect act on the conduction electrons as a contact potential that differs between its configurations by $2V_{\mathrm{d}}$, $\pm V_{\mathrm{d}}$ over a volume $\Omega$ around the defect, as in \cref{app:E.10}. With field operators $\psi_\sigma$ the coupling is $\sigma_zV_{\mathrm{d}}\Omega\sum_\sigma\psi_\sigma^\dagger(0)\psi_\sigma(0)$, so the bath coordinate is $X = 2V_{\mathrm{d}}\Omega\sum_\sigma\psi_\sigma^\dagger(0)\psi_\sigma(0)$, measured from its mean. For a bath of oscillators, \cref{eq:C.20-J} gives at $T = 0$ the spectral function $\int\dd t\,\ee^{\ii\omega t}\langle X(t)X(0)\rangle = 2J_{\mathrm{s}}(\omega)$ for $\omega > 0$, since each mode contributes $2\pi\lambda_k^2\delta(\omega - \omega_k)$. For the electrons the same spectral function, to second order in $V_{\mathrm{d}}$, is a sum over electron--hole pairs. With $\psi_\sigma(0) = \Vol^{-1/2}\sum_{\vect{k}}c_{\vect{k}\sigma}$, a pair with the electron at $\varepsilon'$ above the Fermi level and the hole at $\varepsilon$ below it has the matrix element $2V_{\mathrm{d}}\Omega/\Vol$, and the sums over states become $\Vol\mathcal{N}_{\mathrm{F}}\dd\varepsilon$ each:
\begin{equation*}
  2J_{\mathrm{s}}^{\mathrm{el}}(\omega) = 2\pi\sum_\sigma(2V_{\mathrm{d}}\Omega)^2\mathcal{N}_{\mathrm{F}}^2\int_{-\infty}^0\dd\varepsilon\int_0^\infty\dd\varepsilon'\,\delta(\omega - \varepsilon' + \varepsilon)
  = 16\pi(\mathcal{N}_{\mathrm{F}}V_{\mathrm{d}}\Omega)^2\,\omega .
\end{equation*}
Comparing with $J_{\mathrm{s}} = 2\pi\alpha_{\mathrm{ohm}}\omega$ of \cref{eq:C.20-ohmic} gives $\alpha_{\mathrm{ohm}} = 4(\mathcal{N}_{\mathrm{F}}V_{\mathrm{d}}\Omega)^2 = 4K_{\mathrm{b}}$, \cref{eq:16.8-born}, and \cref{eq:16.8-golden} gives \cref{eq:16.8-black}. The factor four is two from the spin sum and two from the difference $2V_{\mathrm{d}}$ between the configurations. The density of states is per spin: written with the density of states of both spins, $n_0 = 2\mathcal{N}_{\mathrm{F}}$, and $K = (n_0V_{\mathrm{d}}\Omega)^2$, the same result reads $\alpha_{\mathrm{ohm}} = K$, and taking $\alpha_{\mathrm{ohm}} = 4K$ with that $K$ overstates the coupling fourfold.

The pair excitations of a Fermi sea are not oscillators, and the Gaussian form fails at higher order in $V_{\mathrm{d}}$. The exact long-time decay of the overlap of the two Fermi seas, the orthogonality catastrophe, matches the ohmic form with the coupling \cref{eq:E.10-ohm}, $\alpha_{\mathrm{ohm}} = (1/2\pi^2)\sum_\sigma\sum_\nu\tilde\delta_\nu^2$, with $\tilde\delta_\nu$ the differences of the scattering phases at the Fermi level (\cref{app:E.10}). For a contact potential $V$ the scattering in the $s$ channel is the geometric series $t = V\Omega + V\Omega\,G(0)\,V\Omega + \dots = V\Omega/(1 - V\Omega\,G(0))$, with $G(0) = -\ii\pi\mathcal{N}_{\mathrm{F}}$ the imaginary part of the local Green's function at the Fermi level; the real part, which depends on the band away from the Fermi level, renormalizes $V$ and is absorbed into it. Unitarity, $S = 1 - 2\pi\ii\mathcal{N}_{\mathrm{F}}t = \ee^{2\ii\delta}$, then gives
\begin{equation*}
  \ee^{2\ii\delta} = \frac{1 - \ii\pi\mathcal{N}_{\mathrm{F}}V\Omega}{1 + \ii\pi\mathcal{N}_{\mathrm{F}}V\Omega} , \qquad \tan\delta = -\pi\mathcal{N}_{\mathrm{F}}V\Omega .
\end{equation*}
The two configurations, $V = \pm V_{\mathrm{d}}$, have phases $\mp\arctan(\pi\mathcal{N}_{\mathrm{F}}V_{\mathrm{d}}\Omega)$, so $\tilde\delta = 2\arctan(\pi K_{\mathrm{b}}^{1/2})$ for each spin, and
\begin{equation}
  \alpha_{\mathrm{ohm}} = \frac{4}{\pi^2}\arctan^2\bigl(\pi K_{\mathrm{b}}^{1/2}\bigr) = 4K_{\mathrm{b}}\Bigl(1 - \frac{2\pi^2}{3}K_{\mathrm{b}} + O(K_{\mathrm{b}}^2)\Bigr) .
  \label{eq:D.42-alpha}
\end{equation}
The Born limit is the first term. The coupling saturates at $\alpha_{\mathrm{ohm}} = 1$ when $\tilde\delta = \pi$, the largest a single channel per spin can give; larger values need several channels, as for a hopping scatterer, whose displacement mixes the angular momenta about any fixed center (\cref{app:E.22}).

\subsection*{The superconducting factor}

In a superconductor the contact coupling scatters Bogoliubov quasiparticles. For $\Delta_{\mathrm{tl}} < 2\Delta_{\mathrm{BCS}}$ it cannot create a pair, and the only processes are the scattering of a quasiparticle of energy $E$ to $E' = E \pm \Delta_{\mathrm{tl}}$. The potential is even under time reversal, so the matrix element between $\vect{k}$ and $\vect{k}'$ carries the case-I factor $(uu' - vv')^2 = \tfrac12[1 + (\xi\xi' - \Delta_{\mathrm{BCS}}^2)/EE']$ of \cref{eq:7.2-cohE}. Each energy $E > \Delta_{\mathrm{BCS}}$ comes from the two values $\xi = \pm(E^2 - \Delta_{\mathrm{BCS}}^2)^{1/2}$, and summing the factor over the four sign combinations of $\xi$ and $\xi'$ cancels the term in $\xi\xi'$ and leaves $2(1 - \Delta_{\mathrm{BCS}}^2/EE')$. Changing variables from $\xi$ to $E$ brings the density-of-states factor $\rho_{\mathrm{s}}(E) = E/(E^2 - \Delta_{\mathrm{BCS}}^2)^{1/2}$ for each quasiparticle. In the normal state the same golden rule, with constant matrix elements, gives $\int_{-\infty}^\infty\dd\xi\,[f(\xi)(1 - f(\xi + \Delta_{\mathrm{tl}})) + f(\xi + \Delta_{\mathrm{tl}})(1 - f(\xi))] = \Delta_{\mathrm{tl}}\coth(\Delta_{\mathrm{tl}}/2T)$ for emission plus absorption, with $f$ the Fermi function. The ratio of the two rates is therefore
\begin{equation}
  \begin{split}
  \frac{\varGamma_1^{\mathrm{s}}}{\varGamma_1^{\mathrm{el}}} = {}&\frac{2}{\Delta_{\mathrm{tl}}\coth(\Delta_{\mathrm{tl}}/2T)}\int_{\Delta_{\mathrm{BCS}}}^\infty\dd E\;
  \frac{EE' - \Delta_{\mathrm{BCS}}^2}{(E^2 - \Delta_{\mathrm{BCS}}^2)^{1/2}(E'^2 - \Delta_{\mathrm{BCS}}^2)^{1/2}} \\
  &\times\bigl[f(E)(1 - f(E')) + f(E')(1 - f(E))\bigr] ,
  \end{split}
  \label{eq:D.42-Rs}
\end{equation}
with $E' = E + \Delta_{\mathrm{tl}}$. The singularity at $E = \Delta_{\mathrm{BCS}}$ is integrable and is removed by $E = \Delta_{\mathrm{BCS}} + x^2$. \Cref{fig:16-tlsrates,fig:16-limit-tlsrel} use this integral with the weak-coupling gap $\Delta_{\mathrm{BCS}}(T)$.

Two limits are closed. As $\Delta_{\mathrm{tl}} \to 0$ the coherence fraction tends to one, the occupation bracket to $2f(E)(1 - f(E)) = -2T\,\partial f/\partial E$, and $\Delta_{\mathrm{tl}}\coth(\Delta_{\mathrm{tl}}/2T) \to 2T$, so that the ratio becomes $2\int_{\Delta_{\mathrm{BCS}}}^\infty\dd E\,(-\partial f/\partial E) = 2f(\Delta_{\mathrm{BCS}}) = 2/(1 + \ee^{\Delta_{\mathrm{BCS}}/T})$, the factor of Meißner et al.~\cite[Eq.~(9)]{meissner2021acoustic}. For $T$ and $\Delta_{\mathrm{tl}}$ both small compared with $\Delta_{\mathrm{BCS}}$, but in any ratio to each other, write $E = \Delta_{\mathrm{BCS}} + y$ with $y \sim T$. To leading order in $y/\Delta_{\mathrm{BCS}}$ and $\Delta_{\mathrm{tl}}/\Delta_{\mathrm{BCS}}$ the coherence fraction is $(2y + \Delta_{\mathrm{tl}})/2[y(y + \Delta_{\mathrm{tl}})]^{1/2}$, and the occupations are Boltzmann factors, so the bracket is $\ee^{-(\Delta_{\mathrm{BCS}} + y)/T}(1 + \ee^{-\Delta_{\mathrm{tl}}/T})$. With $a = \Delta_{\mathrm{tl}}/2$ and $y = a(\cosh\vartheta - 1)$, so that $y(y + 2a) = a^2\sinh^2\vartheta$ and $y + a = a\cosh\vartheta$,
\begin{equation*}
  \int_0^\infty\dd y\,\ee^{-y/T}\,\frac{y + a}{[y(y + 2a)]^{1/2}} = a\,\ee^{a/T}\int_0^\infty\dd\vartheta\,\ee^{-(a/T)\cosh\vartheta}\cosh\vartheta = a\,\ee^{a/T}K_1(a/T) ,
\end{equation*}
the integral representation of the modified Bessel function $K_1$. With $u = \Delta_{\mathrm{tl}}/2T$, $1 + \ee^{-2u} = 2\ee^{-u}\cosh u$ and $\coth u = \cosh u/\sinh u$, \cref{eq:D.42-Rs} becomes
\begin{equation*}
  \frac{\varGamma_1^{\mathrm{s}}}{\varGamma_1^{\mathrm{el}}} = 2\ee^{-\Delta_{\mathrm{BCS}}/T}\,\sinh u\,K_1(u) ,
\end{equation*}
which is \cref{eq:16.8-sc} to the same order, since $2/(1 + \ee^{\Delta_{\mathrm{BCS}}/T}) = 2\ee^{-\Delta_{\mathrm{BCS}}/T}$ there. From $K_1(u) = 1/u + (u/2)[\ln(u/2) + \gamma_{\mathrm{E}}] - u/4 + O(u^3\ln u)$ the Bessel factor is $1 + (u^2/2)[\ln(u/2) + \gamma_{\mathrm{E}} - \tfrac16] + \dots$, the leading correction to the published factor; for $u \gg 1$ it falls as $(\pi/8u)^{1/2}$. At $T = 0.02\,\Delta_{\mathrm{BCS}}$ and $\Delta_{\mathrm{tl}} = 5T$, where $\Delta_{\mathrm{tl}}/\Delta_{\mathrm{BCS}}$ is no longer small, the integral \cref{eq:D.42-Rs} departs from the closed form by \SI{2}{\percent}.
